\SourceLeaf{334}

\Needspace{4\baselineskip}The following conclusions are clear from this.

I. If these three sines do not all have the same sign, the problem is physically impossible.

In this case the intersection of the two circles described on $b$ and $c$ with angles $\mathfrak{B}, \mathfrak{C}$ will therefore not lie in the prescribed segments, and consequently in neither of them if
\[ \sin(\mathfrak{B}-B), \qquad \sin(\mathfrak{C}-C) \]
have the same sign as each other but
\[ \sin(\mathfrak{A}-A) \]
has the opposite sign. Thus in place of $\mathfrak{B}$ we find $\mathfrak{B}+180^0$, and in place of $\mathfrak{C}$ we find $\mathfrak{C}+180^0$, at the point of intersection.

II. The condition is also sufficient for possibility, as the following proof shows.

1) If the three sines are positive, the intersection of the circles on $b$ and $c$ lies inside the circle through 1, 2, 3.

\textls[120]{Proof.} Suppose it lay outside. There we would have $\mathfrak{B}' = \mathfrak{B}+180^0$ in place of $\mathfrak{B}$ and $\mathfrak{C}' = \mathfrak{C}+180^0$ in place of $\mathfrak{C}$ (so that $\sin(\mathfrak{B}'-B)$ would be negative and likewise $\sin(\mathfrak{C}'-C)$). Yet the angle on $a$ there remains
\[ =360^0-\mathfrak{B}'-\mathfrak{C}'=360^0-\mathfrak{B}-\mathfrak{C}=\mathfrak{A}. \]

2) Similarly, when the three sines are negative, one proves the absurdity of assuming that the intersection of the two auxiliary circles lies inside the main circle.

The point is that $\mathfrak{A}, \mathfrak{B}, \mathfrak{C}$ must necessarily give a sum $=360^0$ (i.e. divisible by $360^0$). The prescribed values have this property of themselves. If they are incompatible, \textls[120]{two} must be changed (by $180^0$) to preserve that property.

If $b, c$ are the two sides on which the auxiliary circles are constructed, the conditions can be stated in the following form.
\[ \sin(\mathfrak{B}-B), \quad \sin(\mathfrak{C}-C) \quad \text{and} \quad \sin(\mathfrak{B}-B+\mathfrak{C}-C) \]
are to have the same sign.
